%%%%%%%%%%%%%%%%%%%%%%%%%%%%%%%%%%%%%%%%%%%%%%%%%
%%%%%%%%%%%% chap: implementation %%%%%%%%%%%%%%%
%%%%%%%%%%%%%%%%%%%%%%%%%%%%%%%%%%%%%%%%%%%%%%%%%

\chapter{Implementation and Functionalities}\label{chapter:implementation}

\section{Page Ingestion and Content Normalization}

The first implemented functionality is page acquisition. NewGEO supports direct page import
and a crawl-oriented ingestion path. In the direct path, a page is imported from markdown,
title, URL, and content type. In the crawl path, the system accepts seeds and optionally raw
HTML or markdown payloads associated with URLs.

The ingestion module normalizes URLs, expands crawl seeds conservatively, and can transform
HTML into readable markdown through an internal parser. Normalized markdown is then stored
in a \texttt{ContentSnapshot}. The snapshot also receives a lightweight deterministic embedding
generated from token hashes. This embedding is not intended to compete with production
semantic encoders, but it is sufficient to validate the retrieval architecture in a fully local
environment.

This functionality matters because external influence begins with what the system decides the
page actually is. If raw content is inconsistent, the benchmark and recommendation stages are
built on unstable input.

\section{Context Packs and Supporting-Document Retrieval}

One of the distinctive functionalities of the application is the context-pack workflow. For a
target page, the user can attach:
\begin{itemize}
    \item a short editorial brief;
    \item voice rules;
    \item supporting read-only documents;
    \item typed constraints.
\end{itemize}

Supporting documents are ranked through cosine similarity between their embeddings and a
query embedding derived from the query-cluster text and context brief. The top results are then
summarized into short contextual notes. In practice, this allows the recommendation engine to
reuse nearby organizational knowledge without editing those supporting sources directly.

The constraint system supports four types:
\begin{itemize}
    \item \texttt{locked\_fact};
    \item \texttt{forbidden\_claim};
    \item \texttt{required\_term};
    \item \texttt{voice\_rule}.
\end{itemize}

This mechanism is central to the thesis because it turns external data from an uncontrolled
influence source into a governed one.

\section{Benchmarking Engine}

Benchmarking is implemented through an engine-connector abstraction. The current prototype
includes three connector modes:
\begin{itemize}
    \item a heuristic comparable benchmark connector;
    \item a non-comparable live spot-check connector;
    \item an API-preview connector for future external integrations.
\end{itemize}

The heuristic benchmark connector computes visibility-oriented metrics from properties of the
candidate page content. These include keyword overlap with the query, early token coverage,
structural signals such as headings and bullet points, authority-like cues, quote friendliness,
and a small deterministic noise term. From these signals, the system derives the stored metrics:
\begin{itemize}
    \item mention rate;
    \item citation share;
    \item answer position;
    \item token share;
    \item quote coverage;
    \item run variance;
    \item overall score.
\end{itemize}

Although these metrics are heuristic, they provide a consistent proxy for visibility. This is
precisely what the prototype needs: a transparent measurement layer that can be executed
offline and later replaced with real connectors without changing the higher-level workflow.

The live spot-check connector reuses the same general scoring path but marks the run as
non-comparable and increases the variance slightly. This enforces an important product rule:
observational checks are useful, but they should not contaminate comparable trend lines.

\section{Recommendation Generation}

Recommendation generation is implemented as a conservative structured rewrite pipeline.
Instead of paraphrasing the whole source aggressively, the generator assembles a draft with the
following typical sections:
\begin{itemize}
    \item an optional top-level title;
    \item a ``Quick answer'' block aligned to the primary query cluster;
    \item a ``Locked facts to preserve'' block;
    \item a ``Related internal context'' block summarizing supporting documents;
    \item the original body under ``Suggested page body''.
\end{itemize}

This choice is deliberate. The goal is to improve surfaceability while minimizing semantic drift.
The rewrite is therefore closer to structured augmentation than to free-form rewriting.

After the draft is built, the system performs two categories of checks:
\begin{enumerate}
    \item explicit constraint evaluation against locked facts, forbidden claims, required terms, and voice rules;
    \item preservation checks for numeric source claims extracted from the original content.
\end{enumerate}

The recommendation bundle also contains a unified markdown diff and a preview score
comparing the baseline page with the projected rewritten page. This enables a reviewer to
assess potential gain before approving publication.

\section{Approval, Export, and Background Processing}

NewGEO does not treat generated text as automatically publishable. A recommendation can be
approved only if the generated bundle exists and no failing constraint checks remain. Once a
recommendation is approved, it can be exported to a markdown file in the configured export
directory.

This separation between generation and export is one of the most important safety features of
the system. It acknowledges that optimized content is high-impact and that external influence
should remain under human supervision.

The worker process extends this workflow by processing queued runs and recommendations in
the background. At the current stage, the worker uses simple polling, but the interface between
queued jobs and execution is already explicit. This is a sound basis for future retry logic,
distributed queues, and richer job observability.

\section{Dashboard Functionalities}

The frontend exposes the implemented workflow through six route-level views:
\begin{itemize}
    \item \textbf{Overview}, which summarizes tracked pages, context packs, runs, recommendations, and score trends;
    \item \textbf{Pages}, which shows the monitored page inventory and attached context-pack status;
    \item \textbf{Queries}, which lists the query clusters that define measured demand;
    \item \textbf{Runs}, which exposes benchmark observations and the distinction between comparable and non-comparable runs;
    \item \textbf{Recommendations}, which presents rationale, confidence, projected lift, and supporting context used;
    \item \textbf{Compare}, which highlights before/after metrics, constraint checks, diffs, and the generated draft.
\end{itemize}

These views matter because the thesis is not only about implementing algorithms. It is about
operationalizing an influence-analysis workflow. The dashboard turns backend records into a
review surface for human decision-making.

\section{Seed Data, Testing, and Verification}

The repository includes a seeded demo project called \textit{Atlas Docs}. The seed creates one
project, one tracked page, one context pack, one query cluster, one recommendation, and one
completed benchmark run. This data serves as a concrete use case for manual exploration and
for regression-style verification of the main workflow.

Backend tests cover:
\begin{itemize}
    \item recommendation generation and locked-fact preservation;
    \item benchmark execution and artifact persistence;
    \item approval and export flow;
    \item crawl-request import behavior;
    \item storage-backend parity between JSON and SQLite;
    \item seed and dashboard route behavior.
\end{itemize}

These tests do not prove the quality of future real-world GEO outcomes, but they do verify
that the currently implemented workflow behaves consistently and that the main architectural
contracts are executable.

\section{Current Limitations}

The current implementation deliberately stops short of a production GEO platform. Several
limitations remain:
\begin{itemize}
    \item benchmark connectors are heuristic rather than live commercial integrations;
    \item embeddings are deterministic hashed vectors rather than semantic encoders;
    \item crawl execution is payload-friendly and conservative rather than a full web crawler;
    \item background execution uses polling instead of a distributed queue;
    \item the frontend is strongest on read-only dashboarding and still lacks full live mutation flows.
\end{itemize}

These limitations do not invalidate the prototype. Instead, they define the exact contribution of
the implementation: it is a functioning first platform layer for studying and governing how
external data can influence LLM responses. That is sufficient for the thesis goal and provides a
clear basis for future work.
